\section{Methodology}
\label{sec:methodology}

\subsection{Sweep Design}

We use a one-parameter-at-a-time (OAT) sweep design~\citep{morris1991}.
For each of the five parameters, we vary that parameter across its range
while holding all other parameters fixed at their DIA default values.

The grid for each parameter is:

\begin{table}[h]
\centering
\caption{Parameter sweep grid. ``Default'' is the DIA baseline value.}
\label{tab:sweep-grid}
\begin{tabular}{llcc}
\toprule
Parameter & Range & Grid points & Default \\
\midrule
$\ufactor$ & $[0.1\%, 5\%]$ & 0.1, 0.5, 1.0, 2.0, 5.0 & 1.0 \\
$\acounty$ & $[1.0, 10.0]$ & 1.0, 2.0, 3.0, 5.0, 10.0 & 3.0 \\
$T$ & $[100, 1{,}000]$ & 100, 200, 400, 600, 1000 & 600 \\
$\aswing$ & $[1.05, 1.20]$ & 1.05, 1.07, 1.10, 1.15, 1.20 & 1.10 \\
$\wvra$ & $[0.2, 0.6]$ & 0.2, 0.3, 0.4, 0.5, 0.6 & 0.4 \\
\bottomrule
\end{tabular}
\end{table}

This gives $5 \times 5 = 25$ non-baseline parameter settings,
plus the one baseline setting, for 26 total pipeline runs.
Each run is a full 50-state sweep using the \texttt{redist states}
command.

\paragraph{Limitation of OAT design.}
One-at-a-time design cannot detect multiplicative interactions between
parameters; the joint sweep described in Section~\ref{sec:joint-sweep}
addresses this limitation.

\subsection{Outcome Metrics}

For each pipeline run, we record:

\begin{enumerate}
\item \textbf{National Democratic seat count} ($D_{\rm nat}$):
      total Democratic seats across all 50 states and 435 districts,
      using 2020 election data.
\item \textbf{State-level partisan gap} ($D_s - 0.5k_s$):
      difference between Democratic seats and the partisan-neutral
      expectation for each state.
\item \textbf{State-average Polsby-Popper} ($\overline{\PP}_s$):
      mean district-level compactness for each state.
\item \textbf{County split count} ($C_s$):
      number of county-district boundaries (lower is better).
\end{enumerate}

\subsection{Statistical Analysis}

For each parameter and each metric, we report:
\begin{itemize}
\item The range of the metric across grid points: $\max - \min$.
\item The standardized sensitivity index: $(\max - \min) / \sigma_0$,
      where $\sigma_0$ is the baseline standard deviation from B.7's
      seed sweep.
\item A binomial test for whether the sweep produces more Democratic
      outcomes at higher vs.\ lower parameter values.
\end{itemize}

The binomial test addresses the concern that a hostile actor could
find a parameter combination that systematically favors one party.

\subsection{The B.0 Bakeoff Baseline}

For context, we compare parameter sensitivity to algorithm sensitivity
as measured in B.0~\citep{b0paper}.
The B.0 bakeoff found a partisan gap of 12.8 percentage points between
unweighted bisection (4D/4R in Wisconsin) and GeoSection (3D/5R).
This is our scale reference: any effect smaller than this is ``large'';
any effect smaller than 1 percentage point is ``negligible''.
